%-----------------------------------LICENSE------------------------------------%
%   This file is free software: you can redistribute it and/or                 %
%   modify it under the terms of the GNU General Public License as             %
%   published by the Free Software Foundation, either version 3 of the         %
%   License, or (at your option) any later version.                            %
%                                                                              %
%   This file is distributed in the hope that it will be useful,               %
%   but WITHOUT ANY WARRANTY; without even the implied warranty of             %
%   MERCHANTABILITY or FITNESS FOR A PARTICULAR PURPOSE.  See the              %
%   GNU General Public License for more details.                               %
%                                                                              %
%   You should have received a copy of the GNU General Public License along    %
%   with this file.  If not, see <https://www.gnu.org/licenses/>.              %
%------------------------------------------------------------------------------%
%   Author: Ryan Maguire                                                       %
%   Date:   2023/12/31                                                         %
%------------------------------------------------------------------------------%
\documentclass{letter}
\usepackage{hyperref}
\hypersetup{colorlinks=true, urlcolor=blue}

\signature{Ryan Maguire}
\address{%
    2-238B\\%
    Mathematics Department\\%
    MIT\\%
    Cambridge, MA USA%
}
\begin{document}
    \begin{letter}{}
        \opening{To whom it may concern,}
            I am writing to apply for the position of physicist at the
            Center for Astrophysics (CfA) working with the atomic and
            molecular physics division.
            Over the past 14 years I have developed a broad and
            interdisciplinary research background in radio astronomy,
            numerical and computational mathematics, optics, and aeronomy.
            In this time I have gained 14 years of Python and shell scripting
            experience, 9 years of experience in Fortran (F77, F90, F95,
            and some F2003 and F2008), a decade of experience in C, and
            have worked in a dozen and a half other languages.%
            \footnote{%
                See \url{https://github.com/ryanmaguire/mitx_mathematics_programming_examples}
                for examples of how I brought programming languages into the
                classroom at MIT (Python, C, Fortran, Pascal, and many more),
                and \url{https://github.com/ryanmaguire/newtonian_black_holes}
                for a hobby project to raytrace a black hole using classical
                physics with the goal of benchmarking the performance of
                different compilers and languages.
            }
            I have learned to combine my knowledge of programming,
            numerical mathematics, and physics to create complex programs
            for analyzing scientific data.
            The best example of this is \texttt{rss\_ringoccs}%
            \footnote{%
                \url{https://github.com/NASA-Planetary-Science/rss_ringoccs}
            },
            a suite of code in C and Python for processing the radio science
            data from the NASA Cassini mission.
            \par
            When I started working at Wellesley college with the radio science
            team, the best methods for obtaining ring profiles of the rings of
            Saturn were limited to just under 1 kilometer resolutions.
            To date, the NASA PDS still only contains 1 kilometer profiles.
            My work in this
            field has allowed us to go much further, pushing the resolution
            down to 10 to 20 meters. We are discovering new density waves in
            the rings and structures that are much finer than anything seen before.
            All of this was made possible by combining mathematics and physics
            together, and developing the codebase to process this.%
            \footnote{
                The bulk of the heavy lifting is in
                \url{https://github.com/ryanmaguire/libtmpl},
                a library I am working on that provides a full implementation
                of \texttt{libm}, the standard C mathematical library,
                and much more (complex numbers, vectors, optics, etc.).
            }
            I believe this combination of skills will greatly benefit
            the research being done at CfA.
        \closing{Sincerely,}
    \end{letter}
\end{document}
